% docs/physics/latex/02_quasielastic_eep.tex - SIMC++ physics note 02: D(e,e'p) and A(e,e'p) quasi-elastic knock-out.
% Standalone PDF docs/physics/pdf/SIMCpp_Physics_02_quasielastic_eep.pdf (copied to src/physics/proton/ and
% src/physics/nuclear/); a chapter of the book SIMC++ Physics Reference. Rules: docs/physics/README.md.
\documentclass{simcnote}
\simccommon
\begin{document}
\simcnotemeta{02}{D(e,e\pr p) and A(e,e\pr p)}{Quasi-elastic proton knock-out: D(e,e\pr p) and A(e,e\pr p)}%
  {Reaction plug-ins \texttt{deuterium\_eep} and \texttt{nucleus\_eep}, physics model \texttt{deforest}, spectral
   functions \texttt{theory} and \texttt{benhar}}
\simcnotemaketitle
\simcabstract{Proton knock-out from deuterium and heavier nuclei in the plane-wave impulse approximation. The six-fold
cross section is the product of an off-shell electron--proton cross section (de Forest's $\sigma_{cc1}$ or
$\sigma_{cc2}$) and a spectral function $S(E_m,P_m)$ that also carries the nuclear transparency as a constant factor.
For deuterium the missing energy is fixed and the proton energy is solved from the kinematics, with a jacobian; for
$A\ge3$ both final energies are generated. Notation, glossary and common physics: note 00.}
\simcnotetoc

% =====================================================================================================================
\section{The reaction in SIMC++}\lbl{sec:reaction}

\subsection{Process and scope}
The reaction is $e + A \to e' + p + (A-1)$ with the proton detected in coincidence with the electron. What is simulated:
the impulse-approximation cross section with a proton spectral function and a constant transparency, and the event
chain of note 00 (radiation on the three legs, energy loss, Coulomb correction of the electron for $Z>1$, both
spectrometers). What is not: final-state interactions other than the constant attenuation, meson-exchange currents,
neutron knock-out (the hadron is always a proton, with proton form factors), correlated two-nucleon emission beyond what
the spectral function contains, and the Coulomb distortion of the proton.

\subsection{Components}
\begin{simctable}{footnotesize}{L{0.15\linewidth}L{0.2\linewidth}L{0.17\linewidth}L{0.23\linewidth}L{0.25\linewidth}}
\caption{The SIMC++ components of D(e,e\pr p) and A(e,e\pr p) (one component per row).}\lbl{tab:components}\\
\toprule\rowcolor{simchead}\simcth{role} & \simcth{interface} & \simcth{registered name} & \simcth{class} & \simcth{file}\\\midrule\endfirsthead
\toprule\rowcolor{simchead}\simcth{role} & \simcth{interface} & \simcth{registered name} & \simcth{class} & \simcth{file}\\\midrule\endhead
\bottomrule\endlastfoot
reaction, $A=2$ & \code{IReaction} & \code{deuterium\_eep} & \code{NuclearEep} & \file{reactions/NuclearEep.cpp}\\
reaction, $A\ge3$ & \code{IReaction} & \code{nucleus\_eep} & \code{NuclearEep} & \file{reactions/NuclearEep.cpp}\\
event completion & (base class) & -- & \code{reactions::EepReactionBase} & \file{reactions/EepReactionBase.cpp}\\
weight assembly & (function) & -- & \code{reactions::assembleWeight} & \file{reactions/EventChain.cpp}\\
cross section & \code{IPhysicsModel} & \code{deforest} & \code{proton::DeForestModel} & \file{physics/proton/ProtonModels.cpp}\\
cross section, $\sigma_{cc1}$ & \code{IPhysicsModel} & \code{deforest\_cc1} & \code{proton::DeForestCC1} & \file{physics/proton/ProtonModels.cpp}\\
cross section, $\sigma_{cc2}$ & \code{IPhysicsModel} & \code{deforest\_cc2} & \code{proton::DeForestCC2} & \file{physics/proton/ProtonModels.cpp}\\
cross section, on shell & \code{IPhysicsModel} & \code{deforest\_cc1\_onshell} & \code{proton::DeForestCC1OnShell} & \file{physics/proton/ProtonModels.cpp}\\
cross-section kernel & (function) & -- & \code{proton::deForest} & \file{physics/proton/ProtonCrossSections.cpp}\\
spectral function, files & \code{ISpectralFunction} & \code{theory} & \code{nuclear::TheorySpectralFunction} & \file{physics/nuclear/SpectralFunctions.cpp}\\
spectral function, tables & \code{ISpectralFunction} & \code{benhar} & \code{nuclear::BenharSpectralFunction} & \file{physics/nuclear/SpectralFunctions.cpp}\\
table kernels & (functions) & -- & \code{nuclear::SfTable}, \code{TheorySfWeight} & \file{physics/nuclear/SpectralFunctionTables.cpp}\\
theory-file reader & (function) & -- & \code{ReadTheoryFile} & \file{setup/SetupTheory.cpp}\\
beam, target, radiation, arms, outputs & as in note 00 & \code{standard}, \code{hallc\_cryo}, \code{simc\_rad}, \ldots & & \ctabr{tab:plugins}\\
\end{simctable}

\subsection{Selection and override}
With no reaction flag the pipeline selects \code{deuterium\_eep} for $\mathrm{nint}(A)=2$ and \code{nucleus\_eep} for
$\mathrm{nint}(A)\ge3$, both with the model \code{deforest}, whose flavour follows the input \code{deForest\_flag}
(\ctabr{tab:defaults}). The set-up uses $M_h=M_N=M_p$, $M_A=M_d$ and $M_\mathrm{rec}=M_n$ for $A=2$, and for $A\ge3$
$M_A$ = \code{targ\%mass\_amu}$\cdot u$ and $M_\mathrm{rec}$ = \code{targ\%mrec\_amu}$\cdot u$, replaced by $M_A-M_p$
when the two differ by more than 100 MeV. After the theory file has been read, the deuterium case is re-decided as ``one
momentum distribution and $E_\mathrm{Fermi}<1$ MeV''. The spectral function is chosen in the reaction's
\code{init()} from the input (\code{use\_benhar\_sf}) and the target $A$ (\secr{sec:theory}, \secr{sec:benhar}).
JSON overrides (\tabr{tab:json}): the model flavour, and the spectral-function type, file and transparency.

\begin{simctable}{footnotesize}{L{0.36\linewidth}L{0.3\linewidth}L{0.34\linewidth}}
\caption{JSON keys of this reaction and its model (one key per row).}\lbl{tab:json}\\
\toprule\rowcolor{simchead}\simcth{key} & \simcth{effect} & \simcth{example}\\\midrule\endfirsthead
\toprule\rowcolor{simchead}\simcth{key} & \simcth{effect} & \simcth{example}\\\midrule\endhead
\bottomrule\endlastfoot
\code{pipeline.physics} & a fixed flavour of the cross section & \code{"deforest\_cc2"}\\
\code{pipeline.physics.deForest\_flag} & flavour of \code{deforest} ($-1$, 0, 1) & \code{\{"type": "deforest", "deForest\_flag": 1\}}\\
\code{pipeline.reaction.spectral\_function.type} & \code{theory} or \code{benhar} & \code{"benhar"}\\
\code{pipeline.reaction.spectral\_function.file} & the table file (in the data directory) & \code{"benharsf\_4.dat"}\\
\code{pipeline.reaction.spectral\_function.transparency} & $T$ of the Benhar weight & \code{0.36}\\
\end{simctable}

\subsection{Event flow}
\begin{simctable}{footnotesize}{L{0.05\linewidth}L{0.33\linewidth}L{0.62\linewidth}}
\caption{One try of D/A(e,e\pr p) (one step per row; the general chain is \ctabr{tab:chain}).}\lbl{tab:flow}\\
\toprule\rowcolor{simchead}\simcth{\#} & \simcth{C++ call} & \simcth{what is specific to this reaction}\\\midrule\endfirsthead
\toprule\rowcolor{simchead}\simcth{\#} & \simcth{C++ call} & \simcth{what is specific to this reaction}\\\midrule\endhead
\bottomrule\endlastfoot
1 & \code{StandardBeam::generate} & nothing\\
2 & \code{NuclearEep::generate} & both slope pairs; for $A\ge3$ the proton energy then the electron energy, for D the electron energy (\eqr{eq:genheavy})\\
3 & \code{EepReactionBase::completeEv} & D: proton energy and jacobian (\eqr{eq:dquad}, \eqr{eq:djac}); $A\ge3$: $E_m$, $P_m$ (\eqr{eq:heavyem})\\
4 & \code{SimcRadiator::radiate} & photon limits of \eqsr{eq:tail1}{eq:tail23}; $A\ge3$: vertex $E_m$, $P_m$ tested against their limits\\
5 & \code{MonteCarlo::run} & nothing\\
6 & \code{EepReactionBase::completeRecon} & $E_m^\mathrm{rec}$ with the residual kinetic energy (\eqr{eq:recem})\\
7 & \code{NuclearEep::spectralWeight}, \code{assembleWeight} & $S$ from the spectral function, $\sigma$ from \code{deforest}\\
\end{simctable}

% =====================================================================================================================
\section{Kinematics, generated variables and generation volume}\lbl{sec:gen}
\subsection{Generated variables}
\begin{simctable}{small}{L{0.3\linewidth}L{0.33\linewidth}L{0.37\linewidth}}
\caption{Generated variables (one variable per row).}\lbl{tab:generated}\\
\toprule\rowcolor{simchead}\simcth{variable} & \simcth{D(e,e\pr p)} & \simcth{A(e,e\pr p), $A\ge3$}\\\midrule\endfirsthead
\bottomrule\endlastfoot
electron slopes $y'_e$, $x'_e$ & flat & flat\\
proton slopes $y'_p$, $x'_p$ & flat & flat\\
proton energy $E_p$ & solved (\eqr{eq:dquad}) & flat, drawn first\\
electron energy $E'$ & flat & flat, drawn after $E_p$\\
missing energy $E_m$ & fixed (\eqr{eq:dem}) & derived (\eqr{eq:heavyem})\\
missing momentum $P_m$ & derived & derived\\
\end{simctable}
For $A\ge3$ the proton energy is drawn in the part of its range compatible with the summed-energy limits, then the
electron energy in the range left by that proton ($\Sigma$ the limits of $E'+E_p$, \eqr{eq:sumegen}):
\begin{feqa}{reactions/NuclearEep.cpp: NuclearEep::generate}{event.f:289-316}
E_p^\mathrm{lo} &= \max\bigl(E^\mathrm{gen}_{p,\mathrm{min}},\ \Sigma_\mathrm{min}-E'^\mathrm{gen}_\mathrm{max}\bigr),\qquad
E_p^\mathrm{hi} = \min\bigl(E^\mathrm{gen}_{p,\mathrm{max}},\ \Sigma_\mathrm{max}-E'^\mathrm{gen}_\mathrm{min}\bigr),\nonumber\\
E_p &= E_p^\mathrm{lo}+\xi_1(E_p^\mathrm{hi}-E_p^\mathrm{lo}),\qquad
w_\mathrm{gen}\leftarrow w_\mathrm{gen}\,\frac{E_p^\mathrm{hi}-E_p^\mathrm{lo}}{E^\mathrm{gen}_{p,\mathrm{max}}-E^\mathrm{gen}_{p,\mathrm{min}}},\nonumber\\
E'^\mathrm{lo} &= \max\bigl(E'^\mathrm{gen}_\mathrm{min},\ \Sigma_\mathrm{min}-E_p\bigr),\qquad
E'^\mathrm{hi} = \min\bigl(E'^\mathrm{gen}_\mathrm{max},\ \Sigma_\mathrm{max}-E_p\bigr),\nonumber\\
E' &= E'^\mathrm{lo}+\xi_2(E'^\mathrm{hi}-E'^\mathrm{lo}),\qquad
w_\mathrm{gen}\leftarrow w_\mathrm{gen}\,\frac{E'^\mathrm{hi}-E'^\mathrm{lo}}{E'^\mathrm{gen}_\mathrm{max}-E'^\mathrm{gen}_\mathrm{min}},
\lbl{eq:genheavy}
\end{feqa}
with $P_p=\sqrt{E_p^2-M_p^2}$; the try fails when a range is empty. For deuterium only the electron energy is drawn, in
$[\max(E'^\mathrm{gen}_\mathrm{min},\Sigma_\mathrm{min}),\ \min(E'^\mathrm{gen}_\mathrm{max},\Sigma_\mathrm{max})]$ with
the same weight factor. The slopes are converted to angles with \ceqr{eq:physang}.

\subsection{Limits}
\paragraph{Vertex limits.} $P^\mathrm{th}_{i,\mathrm{min}}$, $P^\mathrm{th}_{i,\mathrm{max}}$ denote the edges of the
momentum axis of distribution $i$ of the theory file (\secr{sec:theory}):
\begin{feqa}{setup/SetupLimits.cpp: BuildLimits}{init.f:326-343, 381-386}
\text{D:}\quad & E_m = M_p+M_n-M_d\ (\text{fixed}),\qquad 0\le P_m\le\max(|P^\mathrm{th}_\mathrm{min}|,|P^\mathrm{th}_\mathrm{max}|),\nonumber\\
A\ge3:\quad & 0\le P_m\le\max_i\max(|P^\mathrm{th}_{i,\mathrm{min}}|,|P^\mathrm{th}_{i,\mathrm{max}}|)\ \ (790\ \text{MeV/}c\ \text{with Benhar tables}),\nonumber\\
& E_\mathrm{Fermi}\le E_m\le1000\ \mathrm{MeV},\nonumber\\
& M_\mathrm{rec}^{(\mathrm{min,max})} = M_A - M_p + E_m^{(\mathrm{min,max})},\nonumber\\
& T_\mathrm{rec}^{(\mathrm{min})} = \sqrt{(M^{(\mathrm{max})}_\mathrm{rec})^2+(P_m^{(\mathrm{min})})^2}-M^{(\mathrm{max})}_\mathrm{rec},\nonumber\\
& T_\mathrm{rec}^{(\mathrm{max})} = \sqrt{(M^{(\mathrm{min})}_\mathrm{rec})^2+(P_m^{(\mathrm{max})})^2}-M^{(\mathrm{min})}_\mathrm{rec},
\lbl{eq:vertexedge}
\end{feqa}
where the superscripts (min), (max) denote the vertex limits; with Benhar tables $E_\mathrm{Fermi}=0$.

\paragraph{Missing-energy edge ($A\ge3$).} The input window of $E_m^\mathrm{rec}$ (\code{cuts\%Em\%min/max}) is
widened by the total slop $s_{E_m}$, the sum of the beam, electron and proton contributions:
\begin{feqa}{setup/SetupLimits.cpp: BuildLimits}{init.f:285-302}
s_\mathrm{beam} &= \tfrac12\Delta E_b + (\Delta E_1^\mathrm{max}-\Delta E_1^\mathrm{min}),\qquad
s_e = \tfrac{\mathrm{slop}^\delta_e}{100}P_e^0 + (\Delta E_2^\mathrm{max}-\Delta E_2^\mathrm{min}),\nonumber\\
s_p &= \sqrt{\bigl(r+\tfrac{\mathrm{slop}^\delta_p}{100}P_h^0\bigr)^2+M_p^2}-E_{h,\mathrm{max}} + (\Delta E_3^\mathrm{max}-\Delta E_3^\mathrm{min}),\qquad
r = \sqrt{E_{h,\mathrm{max}}^2-M_p^2},\nonumber\\
s_{E_m} &= s_\mathrm{beam}+s_e+s_p+d_E,\qquad
[E_{m,\mathrm{min}}^\mathrm{edge},E_{m,\mathrm{max}}^\mathrm{edge}] = \bigl[\max(0,\ E^\mathrm{cut}_{m,\mathrm{min}}-s_{E_m}),\ E^\mathrm{cut}_{m,\mathrm{max}}+s_{E_m}\bigr],
\lbl{eq:emedge}
\end{feqa}
($\mathrm{slop}^\delta$ the spectrometer slop in $\delta$; the Coulomb terms are 0 because $V_C$ is the same for every
event). When the input minimum equals the maximum the window is opened to $\pm10^6$ MeV. For deuterium the missing-energy
edge is open ($\pm10^{10}$ MeV).

\paragraph{Photon limit.}
\begin{feqa}{setup/SetupLimits.cpp: BuildLimits}{init.f:411-430}
E_\gamma^\mathrm{tot,max} &= E_b^\mathrm{max} + M_p - E'_\mathrm{min} - E_{h,\mathrm{min}} - E_m^{(\mathrm{min})} - T_\mathrm{rec}^{(\mathrm{min})},\nonumber\\
A\ge3:\quad E_\gamma^\mathrm{tot,max} &\leftarrow \min\bigl(E_\gamma^\mathrm{tot,max},\ E^\mathrm{edge}_{m,\mathrm{max}}-E_m^{(\mathrm{min})}
 + T_\mathrm{rec}^{(\mathrm{max})}-T_\mathrm{rec}^{(\mathrm{min})}\bigr),\nonumber\\
& E_m^{(\mathrm{min})}\leftarrow\max\bigl(E_m^{(\mathrm{min})},\ E^\mathrm{edge}_{m,\mathrm{min}}-E_\gamma^\mathrm{tot,max}\bigr),\qquad
E_m^{(\mathrm{max})}\leftarrow\min\bigl(E_m^{(\mathrm{max})},\ E^\mathrm{edge}_{m,\mathrm{max}}\bigr),
\lbl{eq:egammatot}
\end{feqa}
(0 without radiation; \code{Egamma\_gen\_max} when set); every enabled tail has $E^\mathrm{max}_{\gamma i}=E^\mathrm{tot,max}_\gamma$.

\paragraph{Summed generated energy.} $\Sigma$ limits $E'$ for D and $E'+E_p$ for $A\ge3$:
\begin{feqa}{setup/SetupLimits.cpp: BuildLimits}{init.f:442-469}
A\ge3:\quad \Sigma_\mathrm{max} &= \min\bigl(E_b^\mathrm{max}+M_p-T_\mathrm{rec}^{(\mathrm{min})}-E_m^{(\mathrm{min})},\
 E'_\mathrm{max}+E_{h,\mathrm{max}}+E_\gamma^\mathrm{tot,max}\bigr),\nonumber\\
\Sigma_\mathrm{min} &= \max\bigl(E_b^\mathrm{min}+M_p-T_\mathrm{rec}^{(\mathrm{max})}-E_m^{(\mathrm{max})}-E_{\gamma1}^\mathrm{max},\
 E'_\mathrm{min}+E_{h,\mathrm{min}}\bigr),\nonumber\\
\text{D:}\quad \Sigma_\mathrm{max} &= \min\bigl(E_b^\mathrm{max}+M_p-E_{h,\mathrm{min}}-E_m^{(\mathrm{min})}-T_\mathrm{rec}^{(\mathrm{min})},\
 E'_\mathrm{max}+E_{\gamma2}^\mathrm{max}\bigr),\nonumber\\
\Sigma_\mathrm{min} &= \max\bigl(E_b^\mathrm{min}+M_p-E_{h,\mathrm{max}}-E_m^{(\mathrm{max})}-T_\mathrm{rec}^{(\mathrm{max})}-E_\gamma^\mathrm{tot,max},\
 E'_\mathrm{min}\bigr),\nonumber\\
\Sigma &\leftarrow\bigl[\max(0,\Sigma_\mathrm{min}-d_E),\ \Sigma_\mathrm{max}+d_E\bigr],
\lbl{eq:sumegen}
\end{feqa}
with the vertex limits after \eqr{eq:egammatot}. The generation ranges of the energies are
\begin{feqa}{setup/SetupLimits.cpp: BuildLimits}{init.f:477-510}
\text{D:}\quad & [E'^\mathrm{gen}_\mathrm{min},E'^\mathrm{gen}_\mathrm{max}] = \Sigma\cap[E'_\mathrm{min},\ E'_\mathrm{max}+E^\mathrm{max}_{\gamma2}],\nonumber\\
A\ge3:\quad & [E'^\mathrm{gen}_\mathrm{min},E'^\mathrm{gen}_\mathrm{max}] = [\Sigma_\mathrm{min}-E_{h,\mathrm{max}}-E^\mathrm{max}_{\gamma3},\ \Sigma_\mathrm{max}-E_{h,\mathrm{min}}]
 \cap[E'_\mathrm{min},\ E'_\mathrm{max}+E^\mathrm{max}_{\gamma2}],\nonumber\\
& [E^\mathrm{gen}_{p,\mathrm{min}},E^\mathrm{gen}_{p,\mathrm{max}}] = [\Sigma_\mathrm{min}-E'_\mathrm{max}-E^\mathrm{max}_{\gamma2},\ \Sigma_\mathrm{max}-E'_\mathrm{min}]
 \cap[E_{h,\mathrm{min}},\ E_{h,\mathrm{max}}+E^\mathrm{max}_{\gamma3}].
\lbl{eq:genranges}
\end{feqa}

\subsection{Generation volume}
\begin{feq}{reactions/EepReactionBase.cpp: EepReactionBase::genVolume}{simc.f:390-401}
V_\mathrm{gen}^\mathrm{D} = \Delta\Omega_e\,\Delta\Omega_p\,\Delta E'\ \ [\mathrm{sr^2\,MeV}],\qquad
V_\mathrm{gen}^{A\ge3} = \Delta\Omega_e\,\Delta\Omega_p\,\Delta E'\,\Delta E_p\ \ [\mathrm{sr^2\,MeV^2}],
\lbl{eq:genvol}
\end{feq}
with the full ranges $\Delta E = E^\mathrm{gen}_\mathrm{max}-E^\mathrm{gen}_\mathrm{min}$; the event-by-event narrowing is in
$w_\mathrm{gen}$. Without radiation an $A\ge3$ event is kept only if its vertex $E_m$ and $P_m$ lie inside the vertex
limits; with radiation the same test is made after tail 1.

% =====================================================================================================================
\section{Completing the event}\lbl{sec:complete}
\subsection{D(e,e\pr p): the proton energy}
The missing energy is the deuteron binding energy and the residual system is a neutron:
\begin{feq}{reactions/EepReactionBase.cpp: EepReactionBase::completeEv}{event.f:567-568}
E_m = M_p + M_n - M_d\ (=2.2246\ \mathrm{MeV}),\qquad M_\mathrm{rec} = M_d-M_p+E_m = M_n.
\lbl{eq:dem}
\end{feq}
Energy conservation $\nu + M_d = E_p + \sqrt{M_\mathrm{rec}^2 + |\vec q-\vec p|^2}$ at fixed directions gives a quadratic
equation for $E_p$. With $\cos\gamma = \hat q\cdot\hat p$:
\begin{feqa}{reactions/EepReactionBase.cpp: EepReactionBase::completeEv}{event.f:570-579}
a &= -|\vec q|\cos\gamma,\qquad b = |\vec q|^2,\qquad c = \nu+M_d,\qquad t = c^2-b+M_p^2-M_\mathrm{rec}^2,\nonumber\\
Q_A &= 4(a^2-c^2),\qquad Q_B = 4ct,\qquad Q_C = -4a^2M_p^2-t^2,\qquad
E_p = \frac{-Q_B-\sqrt{Q_B^2-4Q_AQ_C}}{2Q_A},
\lbl{eq:dquad}
\end{feqa}
the squared form of $t-2cE_p = 2aP_p$. Since $c^2-a^2\ge c^2-|\vec q|^2 = M_d^2+2M_d\nu-Q^2>0$, $Q_A<0$ and the chosen
root is the larger one. The try fails if the discriminant is negative or $E_p\le M_p$; the sign condition lost in the
squaring is not checked. When the other root lies between the proton energy edges, a warning is printed for the first
5000 tries. Then $P_p$, $\delta_p$ and
\begin{feq}{reactions/EepReactionBase.cpp: EepReactionBase::completeEv}{event.f:939-941}
P_m = P_\mathrm{miss},\qquad T_\mathrm{rec} = \sqrt{M_\mathrm{rec}^2+P_m^2}-M_\mathrm{rec}.
\lbl{eq:dtrec}
\end{feq}

\subsection{A(e,e\pr p): the missing energy}
Both energies are generated, so the missing quantities follow directly from \ceqr{eq:pm}:
\begin{feq}{reactions/EepReactionBase.cpp: EepReactionBase::completeEv}{event.f:942-946}
P_m = P_\mathrm{miss},\qquad M^*_\mathrm{rec} = \sqrt{E_\mathrm{miss}^2-P_m^2},\qquad E_m = M_p + M^*_\mathrm{rec} - M_A,\qquad
T_\mathrm{rec} = \sqrt{M^{*2}_\mathrm{rec}+P_m^2}-M^*_\mathrm{rec},
\lbl{eq:heavyem}
\end{feq}
i.e.\ $E_m$ is defined through the invariant mass $M^*_\mathrm{rec}$ of the unobserved system. When
$E_\mathrm{miss}^2<P_m^2$ the Fortran produces NaN (which then passes the vertex tests); the default mode of SIMC++ fails
the try. The vertex $P_{m,\perp}$ is computed by the Fortran before $P_m$ is assigned, i.e.\ with the previous try's
value; the default mode uses the current value (no output depends on it).

\subsection{Reconstruction}
The reconstruction uses the nominal residual mass and the recoil kinetic energy:
\begin{feq}{reactions/EepReactionBase.cpp: EepReactionBase::completeRecon}{event.f:1350-1352}
T_\mathrm{rec}^\mathrm{rec} = \sqrt{(P_m^\mathrm{rec})^2+M_\mathrm{rec}^2}-M_\mathrm{rec},\qquad
E_m^\mathrm{rec} = \nu^\mathrm{rec}+M_p-E_p^\mathrm{rec}-T^\mathrm{rec}_\mathrm{rec},
\lbl{eq:recem}
\end{feq}
that is $\nu - T_p - T_\mathrm{rec}$, not the invariant-mass definition of \eqr{eq:heavyem}; the two agree when
$M^*_\mathrm{rec}=M_\mathrm{rec}$. With hard cuts, events with $E_m^\mathrm{rec}$ above \code{cuts\%Em\%max} are not
successes.

% =====================================================================================================================
\section{Jacobians}\lbl{sec:jac}
\paragraph{D(e,e\pr p).} Integrating the six-fold cross section over the $E_m$ delta function at fixed $E'$ and
angles gives the factor $|\mathrm{d}E_p/\mathrm{d}E_m|$:
\begin{feq}{reactions/EepReactionBase.cpp: EepReactionBase::completeEv}{event.f:603-605}
J_E = \Bigl|\frac{t\,(c-E_p)+2cE_p(E_p-c)}{2(a^2-c^2)E_p+ct}\Bigr| = \frac{P_p\,E_\mathrm{rec}}{|c\,P_p+a\,E_p|},\qquad E_\mathrm{rec}=c-E_p,
\lbl{eq:djac}
\end{feq}
the second form following from $t-2cE_p=2aP_p$. It is $|\mathrm{d}E_p/\mathrm{d}E_m|$ for
$E_m = \nu+M_d-E_p-\sqrt{M_\mathrm{rec}^2+P_m^2}+\mathrm{const}$, the residual-kinetic-energy definition of \eqr{eq:recem}.
The full jacobian is
\begin{feq}{reactions/EepReactionBase.cpp: EepReactionBase::completeEv}{event.f:603-605, 1015-1016}
J^\mathrm{D} = J_E\,\bigl(1+x_e'^2+y_e'^2\bigr)^{-3/2},
\lbl{eq:djactot}
\end{feq}
without the proton solid-angle factor, although the proton slopes are generated flat and $\Delta\Omega_p$ is part of
$V_\mathrm{gen}$; the factor $1/r_p^3$ would range from 1 at the centre to about 0.985 at the corners of a
$\pm60\times\pm80$ mrad range. Both modes keep the Fortran (an open question).

\paragraph{A(e,e\pr p).}
\begin{feq}{reactions/EepReactionBase.cpp: EepReactionBase::completeEv}{event.f:1015-1025}
J^{A} = \bigl(1+x_e'^2+y_e'^2\bigr)^{-3/2}\bigl(1+x_p'^2+y_p'^2\bigr)^{-3/2}.
\lbl{eq:hjac}
\end{feq}

% =====================================================================================================================
\section{Cross-section model}\lbl{sec:xsec}
\subsection{What the reaction requires from a model and a spectral function}
\code{EepReactionBase::completeMain} asks the reaction's \code{ISpectralFunction} for $S$ (\code{weight(vertex)}, from
the vertex $E_m$ and $P_m$) and the run's \code{IPhysicsModel} for $\sigma$ (\code{sigcc(vertex)}), then forms
\begin{feq}{reactions/EventChain.cpp: assembleWeight}{event.f:1536-1538, 1561-1562}
w = S\cdot J\cdot w_\mathrm{gen}\cdot\sigma\cdot\bigl(1+V_C/E_b\bigr)^2\quad(\text{Coulomb factor with \code{using\_Coulomb}; }V_C=0\text{ for D}).
\lbl{eq:weight}
\end{feq}
The code factorises $\mathrm{d}^6\sigma = K\,S\,\sigma_{ep}$: the model returns $\mathrm{d}^6\sigma/S = K\sigma_{ep}
= E'P_p\sigma_\mathrm{Mott}\sum a_iW_i$, in $\mu$b\,MeV$^2$/sr$^2$ (code comment \fsrc{physics\_proton.f:51-62}). A
spectral function must return a density per proton momentum in MeV$^{-3}$ for D (the jacobian $J_E$ is dimensionless)
and per momentum and missing energy in MeV$^{-4}$ for $A\ge3$; $w$ is then in $\mu$b/(MeV\,sr$^2$) for D and
$\mu$b/(MeV$^2$\,sr$^2$) for $A\ge3$, the units of $V_\mathrm{gen}^{-1}$ (\eqr{eq:genvol}). Any model or spectral function
with these units can replace the built-in ones.

\subsection{The model \code{deforest}: de Forest off-shell cross sections}\lbl{sec:deforest}
The flavour is set by \code{deForest\_flag} (0: $\sigma_{cc1}$, 1: $\sigma_{cc2}$, $-1$: $\sigma_{cc1}$ on shell; other
values are refused) or by the registered names of \tabr{tab:components}. The literature reference is de Forest
\citer{deforest}, whose names $\sigma_{cc1}$, $\sigma_{cc2}$ the code uses without citing the paper. The code writes the
four-vector products with the metric $(+,-,-,-)$ and puts them into de Forest's formulas (metric $(-,+,+,+)$) with
reversed signs. The off-shell kinematics ($q^2\equiv|\vec q|^2$, $q_\mu^2=-Q^2$):
\begin{feqa}{physics/proton/ProtonCrossSections.cpp: deForest}{physics\_proton.f:64-72}
\text{flavours }\sigma_{cc1},\sigma_{cc2}:\quad &\bar E = \sqrt{P_m^2+M_p^2},\qquad \bar q^2 = (E_p-\bar E)^2-q^2,\nonumber\\
\text{on shell}:\quad &\bar E = E_p-\nu,\qquad \bar q^2 = q_\mu^2 = -Q^2,
\lbl{eq:offshell}
\end{feqa}
($\bar E$ is the on-shell energy of a proton of momentum $\vec P_m$). The angle $\gamma$ between $\hat p$ and $\hat q$ and
the azimuth $\phi$ of $\hat p$ around $\hat q$ are
\begin{feqa}{physics/proton/ProtonCrossSections.cpp: deForest}{physics\_proton.f:74-88}
\sin\gamma &= \sqrt{\max\bigl(0,\ 1-(\hat q\cdot\hat p)^2\bigr)},\qquad
\cos\phi = \frac{\hat q_y(\hat q_y\hat p_z-\hat q_z\hat p_y)-\hat q_x(\hat q_z\hat p_x-\hat q_x\hat p_z)}{\sin\gamma\,\sqrt{1-\hat q_z^2}},
\lbl{eq:cosphi}
\end{feqa}
($\cos\phi=0$ if $\sin\gamma=0$, limited to $\pm1$). The nucleon form factors are those of \ceqr{eq:bosted}, at $Q^2$
(not at $\bar q^2$):
\begin{feq}{physics/proton/ProtonCrossSections.cpp: deForest}{physics\_proton.f:90-95}
\tau = \frac{Q^2}{4M_p^2},\qquad F_1 = \frac{G_E+\tau G_M}{1+\tau},\qquad \kappa F_2 = \frac{G_M-G_E}{1+\tau},\qquad
K_2 = \frac{\kappa^2F_2^2}{4M_p^2}.
\lbl{eq:f1f2}
\end{feq}
The lepton factors, with $T=\tan^2(\theta_e/2)$:
\begin{feq}{physics/proton/ProtonCrossSections.cpp: deForest}{physics\_proton.f:97-100}
v_C = \Bigl(\frac{Q^2}{q^2}\Bigr)^2,\qquad v_T = T+\frac{Q^2}{2q^2},\qquad v_S = T+\frac{Q^2}{q^2}\cos^2\phi,\qquad
v_I = \frac{Q^2}{q^2}\sqrt{T+\frac{Q^2}{q^2}}\,\cos\phi.
\lbl{eq:leptonic}
\end{feq}

\paragraph{$\sigma_{cc1}$ (flavours 0 and $-1$).}
\begin{feqa}{physics/proton/ProtonCrossSections.cpp: deForest}{physics\_proton.f:102-109, 128}
\Sigma_1 &= (F_1+\kappa F_2)^2,\qquad \Sigma_2 = F_1^2-\bar q^2K_2,\nonumber\\
W_C &= (\bar E+E_p)^2\,\Sigma_2 - q^2\,\Sigma_1,\qquad W_T = -2\bar q^2\,\Sigma_1,\nonumber\\
W_S & = 4P_p^2\sin^2\gamma\;\Sigma_2,\qquad W_I = -4(\bar E+E_p)P_p\sin\gamma\;\Sigma_2,\nonumber\\
\Sigma_W &= \tfrac14\bigl(v_CW_C+v_TW_T+v_SW_S+v_IW_I\bigr).
\lbl{eq:cc1}
\end{feqa}

\paragraph{$\sigma_{cc2}$ (flavour 1).} With the products of the struck-proton four-momentum $\bar p=(\bar E,\vec P_m)$:
\begin{feqa}{physics/proton/ProtonCrossSections.cpp: deForest}{physics\_proton.f:111-115}
\bar p\cdot p' &= \bar EE_p - P_p\,\hat p\cdot\vec P_m,\qquad \bar p\cdot q = \bar E\nu-|\vec q|\,\hat q\cdot\vec P_m,\nonumber\\
p'\cdot q &= E_p\nu-P_p|\vec q|\cos\gamma,\qquad \bar q\cdot q = (E_p-\bar E)\nu - q^2,
\lbl{eq:products}
\end{feqa}
\begin{feqa}{physics/proton/ProtonCrossSections.cpp: deForest}{physics\_proton.f:117-124, 128}
W_C &= \Bigl(\bar EE_p+\tfrac12(-\bar p\cdot p'+M_p^2)\Bigr)F_1^2-\tfrac12q^2F_1\kappa F_2\nonumber\\
&\quad-\Bigl[\bigl(-\bar p\cdot q\,E_p-p'\cdot q\,\bar E\bigr)\nu+\bar EE_p\,q^2_\mu+\bar p\cdot q\;p'\cdot q-\tfrac12(-\bar p\cdot p'-M_p^2)\,q^2\Bigr]K_2,\nonumber\\
W_T &= -(-\bar p\cdot p'+M_p^2)F_1^2-\bar q\cdot q\;F_1\kappa F_2+\bigl(2\,\bar p\cdot q\;p'\cdot q+(-\bar p\cdot p'-M_p^2)\,q^2_\mu\bigr)K_2,\nonumber\\
W_S &= (P_p\sin\gamma)^2\bigl(F_1^2-q^2_\mu K_2\bigr),\nonumber\\
W_I &= P_p\sin\gamma\Bigl(-(\bar E+E_p)F_1^2+\bigl((-\bar p\cdot q-p'\cdot q)\nu+(\bar E+E_p)\,q^2_\mu\bigr)K_2\Bigr),\nonumber\\
\Sigma_W &= v_CW_C+v_TW_T+v_SW_S+v_IW_I\quad(\text{no factor }1/4).
\lbl{eq:cc2}
\end{feqa}

\paragraph{The cross section.}
\begin{feq}{physics/proton/ProtonCrossSections.cpp: deForest}{physics\_proton.f:129}
\sigma_\mathrm{deF} = \sigma_\mathrm{Mott}(E',\theta_e,Q^2)\;\frac{P_p\,\Sigma_W}{\bar E}\qquad[\mu\mathrm{b\,MeV^2/sr^2}],
\lbl{eq:deforest}
\end{feq}
with $\sigma_\mathrm{Mott}$ of \ceqr{eq:mott}. When $|\cos\phi|$ exceeds 1 by more than $10^{-10}$ the default mode prints
a warning (the Fortran tests after limiting the value, so its warning never appears).

\subsection{Spectral function \code{theory}: the theory files}\lbl{sec:theory}
Without \code{use\_benhar\_sf}, and always for deuterium, the spectral function is read from a theory file chosen by
$\mathrm{nint}(A)$: \code{h2.theory} (2), \code{c12.theory} (12), \code{fe56.theory} (56), \code{au197.theory} (197).
Any other $A$ uses \code{c12.theory}; SIMC++ prints a warning and lets the JSON input name another file. A file holds
the number of shells $n_\rho$, an absorption factor $\eta$ and $E_\mathrm{Fermi}$; for each shell the number of protons
$n_i$, the centroid $E_i$, the width $\Gamma_i$ and a normalisation $N_i$; then the momentum distributions $\rho_i(p)$
on an equidistant grid. At start-up
\begin{feqa}{setup/SetupTheory.cpp: ReadTheoryFile}{init.f:856-905}
n_i &\leftarrow \eta\,n_i,\qquad \rho_i(p)\leftarrow\rho_i(p)/N_i,\qquad
P^\mathrm{th}_{i,\mathrm{min}} = p_1-\tfrac12\Delta p,\qquad P^\mathrm{th}_{i,\mathrm{max}} = p_{n}+\tfrac12\Delta p,\nonumber\\
I_i &= \frac{1}{\pi}\Bigl[\frac{\pi}{2}+\arctan\frac{E_i-E_\mathrm{Fermi}}{\Gamma_i/2}\Bigr]\quad(A\ge3),
\lbl{eq:theoryinit}
\end{feqa}
$I_i$ being the fraction of a Lorentzian of full width $\Gamma_i$ centred at $E_i$ that lies above $E_\mathrm{Fermi}$. The
weight is a linear interpolation between bin centres, times a Lorentzian in $E_m$ normalised above $E_\mathrm{Fermi}$:
\begin{feqa}{physics/nuclear/SpectralFunctionTables.cpp: TheorySfWeight}{event.f:1405-1432}
r &= (P_m-P^\mathrm{th}_{i,\mathrm{min}})/\Delta p,\quad k=\mathrm{nint}(r)\ (1\text{ if }0,\ n-1\text{ if }n),\quad
f = r+\tfrac12-k,\nonumber\\
\rho_i(P_m) & = \rho_{i,k}+f\,(\rho_{i,k+1}-\rho_{i,k})\ \ (0\le r\le n),\nonumber\\
\text{D:}\quad S &= \sum_i n_i\,\rho_i(P_m),\nonumber\\
A\ge3:\quad S &= \sum_i n_i\,\rho_i(P_m)\,\frac{1}{\pi I_i}\,\frac{\Gamma_i/2}{(E_m-E_i)^2+(\Gamma_i/2)^2}\,\Theta(E_m-E_\mathrm{Fermi}),
\lbl{eq:theorysf}
\end{feqa}
($\rho_i=0$ outside $0\le r\le n$). The files tabulate $\rho_i$ from $-490$ to $+490$ MeV/$c$ in steps of 10 MeV/$c$; only
$P_m\ge0$ is used. Summing $\rho_i\,4\pi p^2\Delta p$ over the positive half of each block reproduces $N_i$ to 2.5\,\%
or better (our check of the files: 0.9957 for \code{h2.theory}; 0.9632 and 0.9718 for the two shells of
\code{c12.theory}, whose $N_i$ are 0.960 and 0.972), so $\rho_i/N_i$ is a density normalised to about 1 in the tabulated
range and $S$ integrates to $\eta\sum_in_i$, the number of protons times the absorption factor.

\begin{simctable}{footnotesize}{L{0.2\linewidth}L{0.1\linewidth}L{0.12\linewidth}L{0.15\linewidth}L{0.15\linewidth}L{0.14\linewidth}L{0.14\linewidth}}
\caption{Shells of the theory files (one shell per row).}\lbl{tab:theory}\\
\toprule\rowcolor{simchead}\simcth{file} & \simcth{$\eta$} & \simcth{$E_\mathrm{Fermi}$ [MeV]} & \simcth{$n_i$} & \simcth{$E_i$ [MeV]} & \simcth{$\Gamma_i$ [MeV]} & \simcth{$N_i$}\\\midrule\endfirsthead
\toprule\rowcolor{simchead}\simcth{file} & \simcth{$\eta$} & \simcth{$E_\mathrm{Fermi}$ [MeV]} & \simcth{$n_i$} & \simcth{$E_i$ [MeV]} & \simcth{$\Gamma_i$ [MeV]} & \simcth{$N_i$}\\\midrule\endhead
\bottomrule\endlastfoot
\code{h2.theory} & 1.0 & 0.0 & 1 & (2.2, unused) & (0.0, unused) & 1.0\\
\code{c12.theory} & 0.56 & 15.96 & 2 & 38.1 & 20 & 0.960\\
\code{c12.theory} & 0.56 & 15.96 & 4 & 16.2 & 5.0 & 0.972\\
\code{fe56.theory} & 0.36 & 10.18 & 2 & 50 & 18 & 0.947\\
\code{fe56.theory} & 0.36 & 10.18 & 6 & 37 & 14 & 0.967\\
\code{fe56.theory} & 0.36 & 10.18 & 10 & 23 & 6 & 0.972\\
\code{fe56.theory} & 0.36 & 10.18 & 2 & 16.7 & 3 & 0.946\\
\code{fe56.theory} & 0.36 & 10.18 & 6 & 11.3 & 1.0 & 0.975\\
\code{au197.theory} & 0.22 & 5.78 & 2 & 46 & 19 & 0.903\\
\code{au197.theory} & 0.22 & 5.78 & 6 & 41 & 17 & 0.950\\
\code{au197.theory} & 0.22 & 5.78 & 10 & 32 & 14 & 0.960\\
\code{au197.theory} & 0.22 & 5.78 & 2 & 28 & 12 & 0.942\\
\code{au197.theory} & 0.22 & 5.78 & 14 & 22.9 & 9.4 & 0.966\\
\code{au197.theory} & 0.22 & 5.78 & 6 & 17.2 & 7.9 & 0.954\\
\code{au197.theory} & 0.22 & 5.78 & 18 & 12.5 & 6.0 & 0.971\\
\code{au197.theory} & 0.22 & 5.78 & 12 & 7.7 & 4.0 & 0.974\\
\code{au197.theory} & 0.22 & 5.78 & 6 & 8.3 & 3.7 & 0.958\\
\code{au197.theory} & 0.22 & 5.78 & 3 & 6.2 & 1.5 & 0.957\\
\end{simctable}
The absorption factor $\eta$ is the nuclear transparency of this option: a constant of the file, not an input.

\subsection{Spectral function \code{benhar}: the Benhar tables}\lbl{sec:benhar}
With \code{use\_benhar\_sf}=1 and $A\ge3$ a table \code{benharsf\_A.dat} is read: each line holds $P_m$, $E_m$, the
proton and the neutron bin probabilities and the bin widths $\Delta P_m$, $\Delta E_m$; the proton column is used and
normalised to a total of 1. Any $A$ without its own table uses the C-12 table.

\begin{simctable}{footnotesize}{L{0.25\linewidth}L{0.08\linewidth}L{0.14\linewidth}L{0.13\linewidth}L{0.13\linewidth}L{0.27\linewidth}}
\caption{Benhar tables (one file per row).}\lbl{tab:benhar}\\
\toprule\rowcolor{simchead}\simcth{file} & \simcth{$A$} & \simcth{grid $P_m\times E_m$} & \simcth{$\Delta P_m$ [MeV/$c$]} & \simcth{$\Delta E_m$ [MeV]} & \simcth{remark}\\\midrule\endfirsthead
\bottomrule\endlastfoot
\code{benharsf\_3mod.dat} & 3 & 50$\times$35 & 19.733 & 1 & ``mod'' version for $^3$He\\
\code{benharsf\_4.dat} & 4 & 50$\times$35 & 19.733 & 1 & one negative node\\
\code{benharsf\_12.dat} & 12 & 40$\times$80 & 20 & 5 & also the fallback\\
\code{benharsf\_56.dat} & 56 & 40$\times$80 & 20 & 5 & \\
\code{benharsf\_64.dat} & 64 & 49$\times$100 & 10 & 2 & proton and neutron columns equal\\
\code{benharsf\_197.dat} & 197 & 40$\times$80 & 20 & 5 & \\
\end{simctable}

The lookup interpolates linearly in $P_m$ between the two bracketing nodes $p_j$, $p_{j+1}$ (constant beyond the first
and last node) and linearly in $E_m$ between $E_k$ and $E_{k+1}$ (extrapolated linearly outside), and sets values below
$10^{-20}$ to 0:
\begin{feqa}{physics/nuclear/SpectralFunctionTables.cpp: SfTable::interpolate}{sf\_lookup.f:118-187}
w_2 &= \frac{P_m-p_j}{p_{j+1}-p_j},\quad w_1 = 1-w_2,\qquad s_{1,2} = w_1\,s(E_{k,k+1},p_j)+w_2\,s(E_{k,k+1},p_{j+1}),\nonumber\\
\mathrm{SF} & = s_1+(E_m-E_k)\frac{s_2-s_1}{E_{k+1}-E_k}.
\lbl{eq:sflookup}
\end{feqa}
The Fortran weight divides the interpolated bin probability by the event's own $P_m^2$ and by fixed widths:
\begin{feq}{physics/nuclear/SpectralFunctions.cpp: BenharSpectralFunction::weight; SpectralFunctionTables.cpp: SfTable::lookupDiff (compat)}{event.f:1403-1404; sf\_lookup.f:92-93}
S = Z\cdot T\cdot\mathrm{SF}_d,\qquad \mathrm{SF}_d^\mathrm{compat} = \frac{\mathrm{SF}(E_m,P_m)}{4\pi\,P_m^2\cdot5\cdot20},
\lbl{eq:benhar}
\end{feq}
with $T$ = \code{transparency} (input). This diverges as $1/P_m^2$ below the first node, is infinite at $P_m=0$ and uses
the wrong widths for the $A$ = 3, 4, 64 tables. The default mode divides each node by the volume of its own bin and
interpolates the densities with the same scheme:
\begin{feq}{physics/nuclear/SpectralFunctionTables.cpp: SfTable::load, SfTable::lookupDiff}{none (replaces sf\_lookup.f:93)}
\mathrm{SF}_d^\mathrm{default} = \text{interpolation (\ref{N02-eq:sflookup}) of } \frac{s(E_k,p_j)}{4\pi\,p_j^2\,\Delta P_{m,j}\,\Delta E_{m,k}}.
\lbl{eq:benharcorr}
\end{feq}
In the Fortran an $E_m$ exactly on the last node reuses the interpolation of the previous call; the default mode uses the
last interval. With Benhar tables $E_\mathrm{Fermi}=0$ and the vertex $P_m$ limit is 790 MeV/$c$ for every file.

\subsection{Fermi motion and transparency}
There is no separate generation of Fermi motion: the initial momentum of the struck proton is $\vec P_m$, fixed by the
generated electron and proton, and its distribution enters only through $S$. The transparency is a constant: $\eta$ of
the theory file (the input \code{transparency} is then ignored) or $T$ with the Benhar tables.

% =====================================================================================================================
\section{Radiation, energy loss and decays}\lbl{sec:rad}
With $\Delta T = \max(T_\mathrm{rec}-T_\mathrm{rec}^{(\mathrm{min})},\ T_\mathrm{rec}^{(\mathrm{max})}-T_\mathrm{rec})$ the
photon limits specific to these reactions are (each then bounded by $E_{\gamma1}^\mathrm{max}$ and widened by $d_E$):
\begin{feqa}{radiation/Radc.cpp: generateRad}{radc.f:229-257, 289-291}
A\ge3:\quad E_{\gamma1} &\in \bigl[E_m-E_m^{(\mathrm{max})}-\Delta T,\ \ E_m-E_m^{(\mathrm{min})}+\Delta T\bigr]
\nonumber\\ &\quad(\text{upper limit also}\le E_m-E^\mathrm{edge}_{m,\mathrm{min}}+\Delta T\text{ when one tail is chosen}),\nonumber\\
\text{D:}\quad E_{\gamma1} &\in \bigl[\Sigma_\mathrm{min}-E',\ \ \min(E_{\gamma1}^\mathrm{max},\ \Sigma_\mathrm{max}-E')\bigr]\nonumber\\
&\quad(\text{lower limit set only when one tail is chosen}),
\lbl{eq:tail1}
\end{feqa}
and for tails 2 and 3 the acceptance edges and the measured missing energy:
\begin{feqa}{radiation/Radc.cpp: generateRad}{radc.f:372-389, 428-446}
E_{\gamma2} &\in\bigl[\max(E'-E'_\mathrm{max},\ E^\mathrm{edge}_{m,\mathrm{min}}-E_m-E_{\gamma1}-\Delta T),\nonumber\\
&\qquad\min(E'-E'_\mathrm{min},\ E^\mathrm{edge}_{m,\mathrm{max}}-E_m-E_{\gamma1}+\Delta T,\ E^\mathrm{tot,max}_\gamma-E_{\gamma1})\bigr],\nonumber\\
E_{\gamma3} &\in\bigl[\max(E_p-E_{h,\mathrm{max}},\ E^\mathrm{edge}_{m,\mathrm{min}}-E_m-E_{\gamma1}-E_{\gamma2}-\Delta T),\nonumber\\
&\qquad\min(E_p-E_{h,\mathrm{min}},\ E^\mathrm{edge}_{m,\mathrm{max}}-E_m-E_{\gamma1}-E_{\gamma2}+\Delta T,\nonumber\\
&\qquad\quad E^\mathrm{tot,max}_\gamma-E_{\gamma1}-E_{\gamma2})\bigr],
\lbl{eq:tail23}
\end{feqa}
each widened by $d_E$ (the tail-2 lower missing-energy bound only when one tail is chosen or the proton does not
radiate); for deuterium the missing-energy edge is open and these terms are void. When all tails radiate at once
(\code{rad\_flag} 3), the deuterium tail-1 lower limit keeps the value of an earlier event. After a tail-1 photon the
event is recomputed at the lower beam energy (D: new proton energy and jacobian; $A\ge3$: new $E_m$, $P_m$, tested
against the vertex limits). The soft and hard factors use elastic electron--proton kinematics (\ceqr{eq:bremkin}): for a
quasi-elastic event $E'$ is below the elastic value and the angle rebuilt there is not the event's angle.

The proton loses energy and scatters as in \csecr{sec:eloss} with $M_h=M_p$; the Coulomb correction applies to the
electron for $Z>1$. Nothing decays.

% =====================================================================================================================
\section{Inputs and outputs}\lbl{sec:io}
\subsection{Inputs}
\begin{simctable}{footnotesize}{L{0.24\linewidth}L{0.12\linewidth}L{0.08\linewidth}L{0.56\linewidth}}
\caption{Inputs specific to D/A(e,e\pr p) (one input per row; every input defaults to 0; the inputs common with
H(e,e\pr p) are in note 01, \S7, and the user manual).}\lbl{tab:inputs}\\
\toprule\rowcolor{simchead}\simcth{input} & \simcth{unit} & \simcth{default} & \simcth{meaning here}\\\midrule\endfirsthead
\toprule\rowcolor{simchead}\simcth{input} & \simcth{unit} & \simcth{default} & \simcth{meaning here}\\\midrule\endhead
\bottomrule\endlastfoot
\code{targ\%A} & -- & 0 & 2 selects \code{deuterium\_eep}, $\ge3$ \code{nucleus\_eep}; chooses the theory or Benhar file\\
\code{targ\%Z} & -- & 0 & proton number (Benhar weight $Z$, Coulomb energy)\\
\code{targ\%mass\_amu} & u & 0 & $M_A$ for $A\ge3$ (forced to $M_d$ for $A=2$)\\
\code{targ\%mrec\_amu} & u & 0 & $M_\mathrm{rec}$ for $A\ge3$ (replaced by $M_A-M_p$ if more than 100 MeV off)\\
\code{deForest\_flag} & code & 0 & 0 $\sigma_{cc1}$, 1 $\sigma_{cc2}$, $-1$ $\sigma_{cc1}$ on shell\\
\code{use\_benhar\_sf} & flag & 0 & Benhar tables instead of theory files ($A\ge3$)\\
\code{transparency} & -- & 0 & $T$ of the Benhar weight (no effect with theory files)\\
\code{cuts\%Em\%min} & MeV & 0 & lower end of the $E_m^\mathrm{rec}$ window ($A\ge3$ edge)\\
\code{cuts\%Em\%max} & MeV & 0 & upper end of the window; with hard cuts also a cut\\
\code{using\_Coulomb} & flag & 0 & Coulomb correction ($Z>1$)\\
\end{simctable}

\subsection{Outputs}
% docs/physics/latex/frag_eep_columns.tex - the (e,e'p) tree layout, \input by the notes 01, 02 and 03.
% Source: upstream results_write.f:28-313 (59ccdf6) and src/output/NtupleLayout.cpp (UpstreamEepColumns,
% EepNtupleColumns). No displayed equations here (the lint of build_physics_pdf.py checks the notes that input it).
The tree \code{h10} has one entry per successful event \fsrc{simc.f:346} and, in the default (upstream 59ccdf6) layout,
53 columns \fsrc{results\_write.f:42-313}; \code{-{}-tree-layout legacy} gives the 46 columns of SIMC 294f662 (without
\code{xbj}, the vertex block \code{qi}\ldots\code{Pmi}, and with \code{thetapq}, \code{phipq} instead).
Columns 1--12 belong to the arm on the beam right (prefix \code{hs}), 13--24 to the arm on the beam left (\code{ss}),
whatever particle it detects \fsrc{results\_write.f:42-129}.

\begin{simctable}{footnotesize}{L{0.07\linewidth}L{0.22\linewidth}L{0.5\linewidth}L{0.15\linewidth}}
\caption{The (e,e$'$p) columns of the tree \code{h10} (default layout).}\label{tab:eepcolumns}\\
\toprule\rowcolor{simchead}\simcth{\#} & \simcth{name} & \simcth{content} & \simcth{unit}\\\midrule\endfirsthead
\toprule\rowcolor{simchead}\simcth{\#} & \simcth{name} & \simcth{content} & \simcth{unit}\\\midrule\endhead
\bottomrule\endlastfoot
1--4 & \code{hsdelta}, \code{hsyptar}, \code{hsxptar}, \code{hsytar} & reconstructed $\delta$, slopes, $y_\mathrm{tar}$ of the right arm & \%, rad, cm\\
5--8 & \code{hsxfp}, \code{hsxpfp}, \code{hsyfp}, \code{hsypfp} & focal-plane track & cm, rad\\
9--11 & \code{hsdeltai}, \code{hsyptari}, \code{hsxptari} & \code{orig} values (true, after radiation) & \%, rad\\
12 & \code{hsytari} & $z_t\sin\theta_0$ & cm\\
13--24 & \code{ss}\ldots & the same for the left arm (24: $-z_t\sin\theta_0$) & \\
25--28 & \code{q}, \code{nu}, \code{Q2}, \code{W} & reconstructed $|\vec q|$, $\nu$, $Q^2$, $W$ ($W$ with $M_p$) & GeV, GeV$^2$\\
29--32 & \code{xbj}, \code{epsilon}, \code{Em}, \code{Pm} & reconstructed $x$, $\varepsilon$, $E_m$, $P_m$ & --, GeV\\
33--35 & \code{qi}, \code{nui}, \code{Q2i} & vertex $|\vec q|$, $\nu$, $Q^2$ & GeV, GeV$^2$\\
36, 38 & \code{Wi}, \code{epsiloni} & \code{main\%W}, \code{main\%epsilon}: always 0 for (e,e$'$p) and (e,e$'$)A (\finding{U-B1}) & \\
37 & \code{xbji} & vertex $x$ & \\
39--40 & \code{Emi}, \code{Pmi} & vertex $E_\mathrm{miss}=\nu+M_A-E_p$ and $P_\mathrm{miss}$ & GeV\\
41 & \code{corrsing} & $P_e^\mathrm{rec} - M_pE_\mathrm{beam}/(2E_\mathrm{beam}\sin^2(\theta_e^\mathrm{rec}/2)+M_p)$ \fsrc{results\_write.f:33-34} & GeV\\
42--44 & \code{Pmx}, \code{Pmy}, \code{Pmz} & $\vec P_m$ in the frame of \fsrc{results\_write.f:35-39}: $-P_{m,x}$,
  $(P_{m,z}\hat q_y-P_{m,y}\hat q_z)/n$, $-(P_{m,y}\hat q_y+P_{m,z}\hat q_z)/n$, $n=\sqrt{\hat q_y^2+\hat q_z^2}$ & GeV\\
45--47 & \code{PmPar}, \code{PmPer}, \code{PmOop} & reconstructed components of $\vec P_m$ & GeV\\
48 & \code{fry} & $-$\code{main\%target\%rastery} & cm\\
49 & \code{radphot} & $E_{\gamma1}+E_{\gamma2}+E_{\gamma3}$ & GeV\\
50 & \code{sigcc} & $\sigma$ at the vertex (with the Coulomb factor) & see \secr{sec:xsec}\\
51 & \code{Weight} & the event weight $w$ & see \secr{sec:xsec}\\
52, 53 & \code{theta\_e}, \code{theta\_p} & reconstructed polar angles & rad\\
\end{simctable}

For these reactions \code{Em} and \code{Pm} are those of \eqr{eq:recem}, \code{Emi} and \code{Pmi} the vertex
$E_\mathrm{miss}$ and $P_\mathrm{miss}$, \code{corrsing} is the hydrogen-elastic diagnostic (without meaning here),
\code{sigcc} is $\sigma_\mathrm{deF}$ (with the Coulomb factor) in $\mu$b\,MeV$^2$/sr$^2$ and \code{Weight} is $w$ of
\eqr{eq:weight}.

% =====================================================================================================================
\section{Validation in SIMC++}\lbl{sec:validation}
\begin{simctable}{footnotesize}{L{0.44\linewidth}L{0.17\linewidth}L{0.39\linewidth}}
\caption{CTest tests that cover D/A(e,e\pr p) (one test per row).}\lbl{tab:tests}\\
\toprule\rowcolor{simchead}\simcth{test} & \simcth{labels} & \simcth{what it checks}\\\midrule\endfirsthead
\toprule\rowcolor{simchead}\simcth{test} & \simcth{labels} & \simcth{what it checks}\\\midrule\endhead
\bottomrule\endlastfoot
\code{l1\_deforest} & l1 & \eqsr{eq:offshell}{eq:deforest} against the Fortran table\\
\code{l1\_theory\_init} & l1 & reading of the theory files (\eqr{eq:theoryinit})\\
\code{l1\_spectral\_files} & l1 & the data files are the Fortran ones\\
\code{l1\_sf\_lookup\_init} & l1 & reading of the Benhar tables\\
\code{l1\_sf\_lookup} & l1 & \eqr{eq:sflookup}\\
\code{l1\_eep\_state} & l1 & set-up values of the reaction\\
\code{l1\_generate\_eepA\_fixedrng} & l1 & \eqr{eq:genheavy} with fixed random numbers\\
\code{l1\_complete\_ev\_eepA} & l1 & the vertex completion\\
\code{l1\_complete\_recon\_ev\_eepA} & l1 & the reconstruction\\
\code{l1\_complete\_main\_eepA} & l1 & $S$, $\sigma$ and the weight\\
\code{l1\_event\_capture\_eepA\_d} & l1, replay & every routine of captured D events\\
\code{l1\_event\_capture\_eepA\_fe56} & l1, replay & every routine of captured Fe-56 events\\
\code{l1\_simcrun\_replay\_eepA\_d} & l1, replay & complete D run from a captured random stream, bitwise\\
\code{l1\_simcrun\_replay\_eepA\_fe56} & l1, replay & complete Fe-56 run, bitwise\\
\code{l1\_simcrun\_ranlux\_up\_eep\_d\_rad\_A} & l1, replay, slow & full D run with RANLUX against the upstream Fortran\\
\code{l1\_simcrun\_ranlux\_up\_eep\_c12\_rad\_A} & l1, replay, slow & full C-12 run against the upstream Fortran\\
\code{l1\_simcrun\_ranlux\_up\_upf1\_eep\_cu64\_bh\_rad\_A} & l1, replay, slow & Cu-64 run with the Benhar table\\
\code{regression\_identity\_eep\_d\_rad} & regression & compat-mode D outputs against the baseline\\
\code{regression\_corrected\_eep\_d\_rad} & regression & default-mode D outputs against the baseline\\
\code{regression\_identity\_eep\_fe56\_rad} & regression & compat-mode Fe-56 outputs\\
\code{regression\_corrected\_eep\_fe56\_rad} & regression & default-mode Fe-56 outputs\\
\code{fortran\_diff\_D033\_F3\_benhar\_sf\_diff} & fortran\_diff & \eqr{eq:benhar} and \eqr{eq:benharcorr} in both modes\\
\code{fortran\_diff\_D033\_F5\_pmper\_stale\_pm} & fortran\_diff & the vertex $P_{m,\perp}$ in both modes\\
\code{fortran\_diff\_D033\_F6\_sf\_lookup\_stale} & fortran\_diff & the last-node lookup in both modes\\
\code{fortran\_diff\_D033\_F7\_heavy\_mrec\_guard} & fortran\_diff & the NaN guard of \eqr{eq:heavyem}\\
\code{fortran\_diff\_D014\_deforest\_cos\_phi\_warning} & fortran\_diff & the $\cos\phi$ warning\\
\end{simctable}
Statistical comparisons with the Fortran (seed 4242, $2\times10^5$ events per run; validation ledger,
\code{project\_control/VALIDATION\_LEDGER.md}, entry of 2026-10-05; FORTRAN\_DIFFERENCES.md \S8): every run passes in
both modes (D, C-12, Fe-56, Au-197, He-4 with the C-12 spectral function, Fe-56 with the Benhar table), with at most 4
failed comparisons of 230. The default-mode \code{wtcontr} is lower by 0.5 to 2.8\,\% because of the corrected hadron
$\delta$ cut of \code{pass\_cuts}, the same shift as measured on the Fortran trees.

% =====================================================================================================================
\section{Fortran origin and Fortran findings}\lbl{sec:origin}
\begin{simctable}{footnotesize}{L{0.42\linewidth}L{0.58\linewidth}}
\caption{Fortran routine $\to$ C++ function (one routine per row; the common routines are in \ctabr{tab:map}).}\lbl{tab:map}\\
\toprule\rowcolor{simchead}\simcth{Fortran (upstream 59ccdf6)} & \simcth{C++ (src/...)}\\\midrule\endfirsthead
\toprule\rowcolor{simchead}\simcth{Fortran (upstream 59ccdf6)} & \simcth{C++ (src/...)}\\\midrule\endhead
\bottomrule\endlastfoot
dbase.f:220-226 (reaction flags) & \csrc{core/RunSetup.cpp: DeriveReactionFlags}\\
dbase.f:272-300, 446-452 (masses) & \csrc{setup/SetupReaction.cpp: BuildReactionAndKinematics}\\
dbase.f:555-627 (spectral-function choice) & \csrc{setup/SetupTheory.cpp: BuildTargetAndTheory}\\
init.f theory\_init & \csrc{setup/SetupTheory.cpp: ReadTheoryFile}\\
init.f limits\_init, l.285-517 & \csrc{setup/SetupLimits.cpp: BuildLimits}\\
event.f generate, l.276-414 & \csrc{reactions/NuclearEep.cpp: NuclearEep::generate}\\
event.f complete\_ev, l.565-605, 914-946 & \csrc{reactions/EepReactionBase.cpp: EepReactionBase::completeEv}\\
event.f complete\_recon\_ev, l.1350-1352 & \csrc{reactions/EepReactionBase.cpp: EepReactionBase::completeRecon}\\
event.f complete\_main, l.1402-1404 & \csrc{physics/nuclear/SpectralFunctions.cpp: BenharSpectralFunction::weight}\\
event.f complete\_main, l.1405-1432 & \csrc{physics/nuclear/SpectralFunctionTables.cpp: TheorySfWeight}\\
sf\_lookup.f sf\_lookup\_init & \csrc{physics/nuclear/SpectralFunctionTables.cpp: SfTable::load}\\
sf\_lookup.f sf\_lookup & \csrc{physics/nuclear/SpectralFunctionTables.cpp: SfTable::interpolate}\\
sf\_lookup.f sf\_lookup\_diff & \csrc{physics/nuclear/SpectralFunctionTables.cpp: SfTable::lookupDiff}\\
physics\_proton.f deForest & \csrc{physics/proton/ProtonCrossSections.cpp: deForest}\\
physics\_proton.f (call of deForest) & \csrc{physics/proton/ProtonModels.cpp: DeForestModel::sigcc}\\
radc.f:229-308, 351-363 & \csrc{radiation/Radc.cpp: generateRad}\\
\end{simctable}

\begin{simcfindings}{footnotesize}{L{0.2\linewidth}L{0.48\linewidth}L{0.32\linewidth}}
\caption{Fortran findings that touch D/A(e,e\pr p) (one finding per row; the common ones are in \ctabr{tab:findings}).}\lbl{tab:findings}\\
\toprule\rowcolor{simchead}\simcth{id} & \simcth{Fortran behaviour} & \simcth{default / compat}\\\midrule\endfirsthead
\toprule\rowcolor{simchead}\simcth{id} & \simcth{Fortran behaviour} & \simcth{default / compat}\\\midrule\endhead
\bottomrule\endlastfoot
\finding{D-033 F3} & Benhar weight divided by the event's $P_m^2$ and fixed widths & node widths / Fortran\\
\finding{D-033 F5} & vertex $P_{m,\perp}$ from the previous try's $P_m$ & current $P_m$ / Fortran\\
\finding{D-033 F6} & last-node lookup uses the previous call & last interval / Fortran\\
\finding{D-033 F7} & NaN $M^*_\mathrm{rec}$ accepted & try fails / Fortran\\
\finding{D-014} & $\cos\phi$ warning never printed & printed / not printed\\
\finding{D-033 F4} & no proton solid-angle jacobian for D & kept in both (open question)\\
\finding{D-033 F1} & other $A$ silently use the C-12 spectral function & kept, with a warning; JSON can name a file\\
\finding{D-019 Egamma\_min(1)} & stale deuterium tail-1 lower limit with \code{rad\_flag} 3 & kept in both\\
\finding{D-033 F8/F9} & Benhar: $E_\mathrm{Fermi}=0$, $P_m$ limit 790 MeV/$c$ & kept in both\\
\finding{D-033 F12} & two-root warning only in the first 5000 tries & kept in both\\
\finding{D-033 doing\_deuterium} & the deuterium case re-decided from the theory file & kept in both\\
\finding{D-033 F6 logsf} & a linear interpolation stored in a variable called \code{logsf} & kept in both\\
\finding{D-033 F8 krel} & \code{krel} never set for (e,e\pr p) & not an output\\
\end{simcfindings}

\begin{simcrefs}
\simcref{deforest} T.~de Forest, Jr., ``Off-shell electron--nucleon cross sections: the impulse approximation'',
  Nucl.\ Phys.\ A \textbf{392}, 232 (1983). Not cited in the code, which uses its names $\sigma_{cc1}$ and $\sigma_{cc2}$.
\end{simcrefs}
\end{document}
